% VI/23 editorial reader bridge — additive, not a canonical problem.
\noindent\textbf{Theory and worked example: tensor products encode compatible maps.}
The fibre product $X\times_S Y$ represents pairs of maps to $X$ and $Y$
whose composites to $S$ agree. For affine schemes this universal property
becomes the tensor-product identity
\[
\operatorname{Spec}B\times_{\operatorname{Spec}A}
\operatorname{Spec}C\ \cong\ \operatorname{Spec}(B\otimes_A C).
\]
It is important that the tensor product is taken over $A$ rather than over
the ground field: the base imposes equations on the two coordinate sets.

For example, let $A=k[t]$, $B=k[u]$, $C=k[v]$ and let $t$ act as $u^2$
on $B$ and as $v^2$ on $C$. Then
\[
B\otimes_A C\cong k[u,v]/(u^2-v^2).
\]
When $\operatorname{char}k\ne2$, this is the reduced union of the two
lines $u=v$ and $u=-v$, meeting at the origin. In characteristic $2$,
$u^2-v^2=(u-v)^2$ and the fibre product is a nonreduced doubled line.
Thus a universal property that looks purely categorical can detect
components and nilpotent structure in a concrete calculation.

\medskip
